\section{The mutation-necessity archive and faithful-comparator challenge}
\label{sec:necessity}

The 66-case study above remains a historical local falsifier and is not
retrospectively promoted. An archived protocol defines the narrower experiment
residual left after nearest-work saturation: whether a high-level scientific
mutation should require lower-level exclusion, a same-task counterfactual
effect, preservation of a protected sibling or invariant, and an exact
dependency-impact binding. Intervention-supported attribution, active diagnosis,
diagnosis-to-valid-recovery admission, counterfactual inclusion-minimal repair,
and selective dependency rollback are treated as donor-owned substrate rather
than as contributions of this paper. In particular, REFLECT \cite{reflect2026},
Causal Agent Replay (arXiv:2606.08275), R2Act (arXiv:2607.04623), DARC
(arXiv:2608.11772), HERALD (arXiv:2608.06012), CausalFlow (arXiv:2605.25338),
CausalRepair (arXiv:2608.10613), and dependency-guided rollback
\cite{dependencyrollback2026} determine the parent and ablation structure. The
round-by-round dispositions and extracted structures are retained in the
assimilation ledgers named in Section~\ref{sec:availability}.

\subsection{Archived design and integrity boundary}

Each run contains 2,882 credential-free mechanical worlds: 480 hidden
formulation/search-universe shifts and 2,402 evidence-only or execution-only
negative controls. The hidden set is balanced across four shift families; the
control size was fixed prospectively for the safety and non-regression estimand.
This before-outcome ordering is a statement in the archived protocol, not an
independently confirmed chronology.
Every arm receives the same public repair and diagnostic menu and the same
four-unit intervention budget. Protected response matrices, family identity,
gold minimal repairs, protected siblings, and dependency-impact sets remain
host-only.

Nine runnable arms include static refusal, immediate outcome-flip and
cost-greedy repair, replay-based attribution in the style of REFLECT and Causal
Agent Replay, active value-of-information recovery, admission-plus-rollback in
the style of DARC and R2Act, a minimal-counterfactual parent in the style of
CausalFlow and HERALD, and the governed policy. Five direct ablations remove
lower-level exclusion, protected-sibling checking, dependency-impact binding,
$K/W/M$ level ordering, or the intervention budget. The registered support
vector also requires zero unnecessary high-level reframes on controls,
whole-grid non-inferiority, zero protected-sibling regressions, zero forbidden
high-level mutations, and mean cost no greater than the maximum strong-parent
mean plus one unit.

The original execution receipt records that it bound a replay envelope rather
than the committed world receipt and failed before any arm executed. A first
additive amendment records the corresponding correction. A later review found
that the cost implementation compared the governed policy with every parent,
although the registered rule names the maximum parent mean. The second additive
amendment states that no arm or terminal had been observed and binds the exact
registered implementation without changing worlds, hypotheses, parents, seeds,
margins, or protocol. These are immutable archived assertions, not independently
confirmed event order. Because three cited provenance commits are unreachable,
the historical prospective order cannot be checked independently; the current
exact replay cannot repair it.

\subsection{Archived primary and disjoint-replication results}

The archived primary run and registered disjoint replication both carry the
original supported terminal. The governed policy attains protected
root-task success 1.0000 on the 480 hidden-shift worlds in each run. Each strong
parent attains 0.4938 in the primary and 0.4833 in replication. The paired
advantage is therefore 0.5063 in the primary and 0.5167 in replication; every
95\% interval excludes zero and exceeds the registered 0.10 practical margin
(Table~\ref{tab:P1-T4}). Exact paired $p$-values pass the three-comparison Holm
family in both runs. These intervals and tests are generator-relative
comparisons over the paired designed worlds. They do not estimate uncertainty
over naturalistic scientific tasks or authorize model-general inference. The
faithful-comparator challenge below further shows that these margins do not
establish comparative mechanism necessity.

\begin{table}[htbp]
\centering
\small
\setlength{\tabcolsep}{4pt}
\caption{Archived V2 hidden-shift contrasts against the three registered non-iterating parents. Rates are protected root-task success over 480 bound hidden-shift worlds per run. The values and paired intervals remain immutable historical results. The later faithful-comparator challenge matched the governed policy on the primary joint criterion, so these margins do not establish comparative mechanism necessity.}
\label{tab:P1-T4}
\begin{tabular}{llrrrr}
\toprule
Run & Registered V2 parent & Governed & Parent & Difference [95\% CI] & Exact $p$ \\
\midrule
Primary & Active value-of-information & 1.0000 & 0.4938 & 0.5063 [0.4625, 0.5521] & $1.41\times10^{-73}$ \\
Primary & Admission + rollback & 1.0000 & 0.4938 & 0.5063 [0.4604, 0.5500] & $1.41\times10^{-73}$ \\
Primary & Minimal counterfactual repair & 1.0000 & 0.4938 & 0.5063 [0.4625, 0.5521] & $1.41\times10^{-73}$ \\
\midrule
Replication & Active value-of-information & 1.0000 & 0.4833 & 0.5167 [0.4708, 0.5604] & $4.42\times10^{-75}$ \\
Replication & Admission + rollback & 1.0000 & 0.4833 & 0.5167 [0.4729, 0.5604] & $4.42\times10^{-75}$ \\
Replication & Minimal counterfactual repair & 1.0000 & 0.4833 & 0.5167 [0.4708, 0.5625] & $4.42\times10^{-75}$ \\
\bottomrule
\end{tabular}
\end{table}


\subsection{Faithful-comparator falsification}

The R4 challenge changed one defect shared by the three registered parents.
Each parent selected only its top-confidence high-level candidate and did not
try another candidate when the first failed admission. The challenge replaced
that single pick with confidence-ordered search while preserving the public
worlds, four-unit budget, probe ladder, branch structure, and admission test.
The unchanged primary anchor arms reproduced their archived rates exactly.

On the 480 primary hidden-shift worlds, the ordered-search active
value-of-information comparator and the governed policy both achieved joint
constraint-respecting success 1.0000. Both had zero forbidden high-level
mutations, and there were zero discordant pairs. The registered falsification
criterion therefore fired. The earlier 0.5063 margin does not distinguish the
governed composition from the absence of candidate iteration in the registered
comparator, and therefore does not establish necessity for the composition.
The archived internal terminal remains unchanged, but its comparative
mechanism-necessity interpretation is withdrawn.

The governed policy used a mean of 1.834 intervention units on the primary
worlds, compared with 2.668 for the matching ordered-search comparator. This is
a bounded descriptive economy result. The challenge did not isolate which
component gives the comparator zero forbidden mutations, because its
non-iterating parent already had that rate before repair.

The replication anchor gate used primary-run expected rates and did not
reproduce the replication parents. Under the registered rule, no comparison is
read from that arm. The replication challenge therefore remains undetermined,
not positive or negative evidence.

On 2,402 controls per run, the governed policy makes zero unnecessary
formulation/search-universe reframes and matches the active
value-of-information parent's protected success at 1.0000; the paired difference
and interval are exactly 0. Across all 2,882 worlds, its protected success is
1.0000. The primary whole-grid advantages over the active value-of-information,
admission-plus-rollback, and minimal-counterfactual parents are 0.0843, 0.3220,
and 0.9056; replication gives 0.0861, 0.3362, and 0.8980. There are zero
protected-sibling regressions and zero forbidden high-level mutations in both
runs. Mean intervention cost is 1.8341 units in the primary and 1.8359 in
replication, below the maximum strong-parent mean of 4.0.

The registered ablations alter the governed implementation, not only its label
(Figure~\ref{fig:P1-7}). Removing the protected-sibling check reduces
hidden-shift protected success to 0.4938 and 0.4833 and introduces whole-grid
sibling-regression rates of 0.0843 and 0.0861. Removing dependency-impact
binding reduces hidden-shift protected success to zero. Removing lower-level
exclusion or $K/W/M$ ordering creates forbidden high-level mutations on
approximately 0.417 and 0.411 of primary worlds, respectively; the same failure
directions recur in replication. These ablations diagnose the archived
implementation. They do not restore comparative necessity after the faithful
comparator matches the joint criterion.

\begin{figure}[htbp]
  \centering
  \includegraphics[width=0.96\linewidth]{P1-7_necessity_replication.pdf}
  \caption{Descriptive rates from the archived primary and disjoint replication. Left: protected success on 480 hidden-shift worlds. Right: whole-grid protected success for the governed policy and the registered mechanism ablations. Confirmatory inference uses the generator-relative paired intervals and gates in Table~\ref{tab:P1-T4}, not visual bar separation; these quantities do not estimate naturalistic task uncertainty.}
  \label{fig:P1-7}
\end{figure}

\subsection{Separate-implementation verification and authorized claim}

A separately implemented standard-library verifier imports neither the
candidate policies, the production scorer and statistics, nor the subject's own
licensing code. It recomputes all 40,348 arm/world score rows per run, the
complete analysis object, and the terminal. Both verification records report
zero score mismatches, zero analysis mismatches, and matching supported
terminals. This is local code-path separation, not external or institutionally
independent verification.

A later clean-source CPU replay reconstructed both world directories and reran
both campaigns plus the separate verifier. For the primary and replication, the
raw results, result records, and verification records were byte-identical to the
archives. Both terminals and all mismatch counts were unchanged. This closes
current deterministic replay from committed source; it does not create external
scorer custody or independent review.

The replay also cannot repair historical chronology. Three commits cited as
original execution-source provenance are unreachable from the repository and
its public remote. The archived bytes and current bound sources are
reproducible, but the claimed original prospective commit order therefore
remains unverifiable. No inference below uses current replay to promote that
historical order.

The authorized claim is now narrower. On the frozen primary world family, the
governed policy implements the registered preservation and dependency-impact
constraints and uses fewer mean intervention units than the matching
ordered-search comparator. The original contrasts against three non-iterating
parents and the direct ablations remain descriptive history. They do not
establish comparative mechanism necessity or general superiority. The R4
replication remains undetermined. Nothing here establishes model-general
performance, open-ended scientific superiority, literature-search adequacy, or
novelty of attribution, minimal repair, action admission, dependency rollback,
or certificate enforcement.
